% =============================================================================
% Chapter 2 / manuscript: goethite stability, dehydroxylation kinetics and the
% controls on 3He release.  Independent sections for weaving into the existing
% architecture; nothing here defines a preamble.
%
% Figures are produced by the scripts in Chapter2/Figures and written to
% generated/:
%   goethite_hematite_phase_diagram.py  -> goethite_hematite_Ph2o_T_hydrated,
%                                          size_basis_sensitivity,
%                                          nucleation_barrier_ratio_vs_T
%   goethite_vacuum_kinetics.py         -> arrhenius_reaction_vs_diffusion,
%                                          vacuum_kinetics_KAF_to_HM,
%                                          first_order_delta_lnk,
%                                          hydrothermal_translation,
%                                          vacuum_kinetics_KAF_to_HM_fits.csv,
%                                          synthetic_degeneracy_test.csv
%   goethite_hydrothermal_tests.py      -> hydrothermal_diffusion_vs_conversion
%
% Citation keys used (placeholders; map to your .bib; verify volume/pages):
%  Thermodynamic data and conventions
%   navrotsky2008   Navrotsky, Mazeina & Majzlan (2008) Science 319, 1635
%   majzlan2003a    Majzlan et al. (2003) Am. Mineral. 88, 846 (S, Cp)
%   majzlan2003b    Majzlan, Grevel & Navrotsky (2003) Am. Mineral. 88, 855
%   voigt1981       Voigt & Will (1981) N. Jb. Miner. Mh. 2, 89
%   berman1985      Berman & Brown (1985) Contrib. Mineral. Petrol. 89, 168
%   berman1988      Berman (1988) J. Petrol. 29, 445 (apparent G convention)
%   helgeson1978    Helgeson et al. (1978) Am. J. Sci. 278-A
%   anderson2005    Anderson (2005) Thermodynamics of Natural Systems, 2nd ed., CUP
%   wagner2002      Wagner & Pruss (2002) J. Phys. Chem. Ref. Data 31, 387 (IAPWS-95)
%   bell2014        Bell et al. (2014) Ind. Eng. Chem. Res. 53, 2498 (CoolProp)
%   robie1995       Robie & Hemingway (1995) USGS Bull. 2131
%  Kinetics and nucleation
%   fechtig1966     Fechtig & Kalbitzer (1966)
%   mampel1940      Mampel (1940) Z. Phys. Chem. A 187, 235
%   gibbs1878       Gibbs (1876-1878) Trans. Conn. Acad. Arts Sci. 3
%   volmer1926      Volmer & Weber (1926) Z. Phys. Chem. 119, 277
%   becker1935      Becker & Doring (1935) Ann. Phys. 416, 719
%   turnbull1950    Turnbull (1950) J. Chem. Phys. 18, 198 (heterogeneous nucleation)
%   eyring1935      Eyring (1935) J. Chem. Phys. 3, 107 (transition-state theory)
%   christian2002   Christian (2002) The Theory of Transformations in Metals and Alloys, Pergamon
%   kashchiev2000   Kashchiev (2000) Nucleation: Basic Theory with Applications, Butterworth-Heinemann
%   deyoreo2003     De Yoreo & Vekilov (2003) Rev. Mineral. Geochem. 54, 57
%   lasaga1998      Lasaga (1998) Kinetic Theory in the Earth Sciences, Princeton
%   porter2009      Porter, Easterling & Sherif (2009) Phase Transformations in Metals and Alloys, 3rd ed., CRC
%  Goethite dehydroxylation
%   limadefaria1963 Lima-de-Faria (1963) Z. Kristallogr. 119, 176
%   watari1979      Watari et al. (1979) J. Solid State Chem. 29, 137
%   watari1983      Watari et al. (1983) J. Solid State Chem. 48, 49
%   cudennec2005    Cudennec & Lecerf (2005) Solid State Sci. 7, 520
%   goss1987        Goss (1987) Mineral. Mag. 51, 437
%   saito2016       Saito et al. (2016) Phys. Chem. Miner. 43, 749
% Items marked TODO need a check before submission.
% =============================================================================


% -----------------------------------------------------------------------------
\section{Thermodynamic setting}
\label{sec:gt_thermo}
% -----------------------------------------------------------------------------

Goethite breaks down to hematite by the dehydroxylation reaction
\begin{equation}
  2\,\alpha\text{-FeOOH} \;\rightleftharpoons\; \alpha\text{-Fe}_2\text{O}_3 + \text{H}_2\text{O}_{(\text{fluid})},
  \label{eq:gt_rxn}
\end{equation}
whose Gibbs energy change, written per mole of hematite formed, is
\begin{equation}
  \Delta G_r(T,P) = G_{\mathrm{hem}}(T,P) + G_{\mathrm{w}}(T,P) - 2\,G_{\mathrm{goe}}(T,P) + \Delta G_{\mathrm{surf}} .
  \label{eq:gt_dGr}
\end{equation}
Goethite is stable where $\Delta G_r>0$ and hematite plus water where $\Delta G_r<0$.

\paragraph{Gibbs energies of the solids.} The Gibbs energy of a phase at $(T,P)$ follows from
integrating $dG = -S\,dT + V\,dP$ from the reference state $(T_r,P_r) = (298.15~\mathrm{K}, 1~\mathrm{bar})$,
first in temperature at $P_r$ and then in pressure at $T$. At $P_r$,
$H(T) = H(T_r) + \int_{T_r}^{T}C_p\,dT$ and $S(T) = S(T_r) + \int_{T_r}^{T}(C_p/T)\,dT$, so
$G(T,P_r) = H(T) - T S(T)$. Adding $\int_{P_r}^{P}V\,dP$, and treating the solids as incompressible
with no thermal expansion, gives
\begin{equation}
  G_i(T,P) = \Delta_f H_i^{0} + \int_{T_r}^{T} C_{p,i}\,dT
  \;-\; T\left[S_i^{0} + \int_{T_r}^{T}\frac{C_{p,i}}{T}\,dT\right]
  \;+\; V_i\,(P-P_r),
  \label{eq:gt_Gsolid}
\end{equation}
where we have set $H(T_r) = \Delta_f H_i^{0}$, the enthalpy of formation from the elements, and
$S(T_r) = S_i^{0}$, the third-law entropy. This is the ``apparent'' Gibbs energy of formation
\citep{berman1988}; see also \citet{helgeson1978} and \citet{anderson2005}. It differs from
the true Gibbs energy of formation,
$\Delta_f G_i = \Delta_f H_i - T(S_i - \sum_e n_{e,i} S_e)$, by the term $T\sum_e n_{e,i}S_e$,
where $n_{e,i}$ is the number of atoms of element $e$ in phase $i$ and $S_e$ the entropy of the
element in its reference state. For any balanced reaction $\sum_i \nu_i n_{e,i} = 0$ for every element,
so these terms cancel in $\Delta G_r$. The incompressible approximation affects $\Delta G_r$
through $\Delta V_s (P-P_r)$ with $\Delta V_s = V_{\mathrm{hem}} - 2V_{\mathrm{goe}} = -11.5$~cm$^3$,
i.e.\ at most a few kJ below 1~GPa; compressibility and thermal expansion change $V_i$ by
$\lesssim1\%$ over the conditions considered and are neglected.

Goethite heat capacity uses the Berman--Brown form \citep{berman1985} refined by
\citet{majzlan2003b},
\begin{equation}
  C_p = k_0 + k_1 T^{-1/2} + k_2 T^{-2} + k_3 T^{-3},
\end{equation}
which integrates to
\begin{align}
  \int_{T_r}^{T} C_p\,dT &= k_0 (T-T_r) + 2k_1\left(T^{1/2}-T_r^{1/2}\right)
    - k_2\left(T^{-1}-T_r^{-1}\right) - \tfrac{1}{2}k_3\left(T^{-2}-T_r^{-2}\right), \\
  \int_{T_r}^{T} \frac{C_p}{T}\,dT &= k_0 \ln\frac{T}{T_r} - 2k_1\left(T^{-1/2}-T_r^{-1/2}\right)
    - \tfrac{1}{2}k_2\left(T^{-2}-T_r^{-2}\right) - \tfrac{1}{3}k_3\left(T^{-3}-T_r^{-3}\right).
\end{align}
Hematite uses a Maier--Kelley form, $C_p = a + bT + cT^{-2}$, with
\begin{align}
  \int_{T_r}^{T} C_p\,dT &= a(T-T_r) + \tfrac{1}{2}b\left(T^2-T_r^2\right) - c\left(T^{-1}-T_r^{-1}\right), \\
  \int_{T_r}^{T} \frac{C_p}{T}\,dT &= a\ln\frac{T}{T_r} + b(T-T_r) - \tfrac{1}{2}c\left(T^{-2}-T_r^{-2}\right).
\end{align}
% TODO: cite the source of the hematite Maier-Kelley coefficients (Kelley-type fit,
% cf. robie1995); Majzlan et al. used Berman (1988). Differences are < 1 J/K/mol at 298 K.

\paragraph{Water.} Water is evaluated with the IAPWS-95 equation of state \citep{wagner2002},
implemented in CoolProp \citep{bell2014}, which is valid for liquid, vapor and supercritical
water to 1~GPa. In the same apparent convention,
\begin{equation}
  G_{\mathrm{w}}(T,P) = \Delta_f H_{\mathrm{w}}^{0} + \left[h(T,P) - h(T_r,P_r)\right]
  - T\left[S_{\mathrm{w}}^{0} + s(T,P) - s(T_r,P_r)\right],
  \label{eq:gt_Gw}
\end{equation}
where $h$ and $s$ are the IAPWS-95 molar enthalpy and entropy and the reference state is liquid
water at $T_r,P_r$. Equation~\ref{eq:gt_Gw} is exact within the equation of state; no ideal-gas
approximation is made anywhere in the calculation. Table~\ref{tab:gt_thermo_data} lists all
input data. We use the enthalpies and entropies of \citet{majzlan2003b} refined by mathematical
programming against the phase-equilibrium brackets of \citet{voigt1981}.

\begin{table}
  \centering
  \caption{Thermodynamic data used for Eq.~\ref{eq:gt_dGr}.}
  \label{tab:gt_thermo_data}
  \begin{tabular}{lcc}
    \hline
    & Goethite ($\alpha$-FeOOH) & Hematite ($\alpha$-Fe$_2$O$_3$) \\
    \hline
    $\Delta_f H^0$ (kJ\,mol$^{-1}$) & $-561.4$ & $-825.8$ \\
    $S^0$ (J\,mol$^{-1}$\,K$^{-1}$) & 59.5 & 87.4 \\
    $V^0$ (cm$^3$\,mol$^{-1}$) & 20.881 & 30.274 \\
    $C_p$ & $k_0=133.8765$, $k_1=-738.5577$, & $a=98.28$, $b=7.782\times10^{-2}$, \\
          & $k_2=-2.177657\times10^{6}$, $k_3=2.165124\times10^{8}$ & $c=-1.485\times10^{6}$ \\
    $C_p(298\ \mathrm{K})$ (J\,mol$^{-1}$\,K$^{-1}$) & 74.8 & 104.8 \\
    $\gamma_{\mathrm{hydrated}}$ (J\,m$^{-2}$) & $0.60\pm0.10$ & $0.75\pm0.16$ \\
    $\gamma_{\mathrm{anhydrous}}$ (J\,m$^{-2}$) & $0.91\pm0.09$ & $1.90\pm0.30$ \\
    \hline
    \multicolumn{3}{l}{\footnotesize Enthalpy, entropy and $C_p$: \citet{majzlan2003a,majzlan2003b}; surface energies: \citet{navrotsky2008}.}\\
    \multicolumn{3}{l}{\footnotesize Water: $\Delta_f H^0_{\mathrm{w}} = -285.83$~kJ\,mol$^{-1}$, $S^0_{\mathrm{w}} = 69.95$~J\,mol$^{-1}$\,K$^{-1}$; IAPWS-95 for $h(T,P)$, $s(T,P)$.}
  \end{tabular}
\end{table}

As a check, at 298.15~K and 1~bar with liquid water, Eqs.~\ref{eq:gt_dGr}--\ref{eq:gt_Gw}
give $\Delta G_r = -0.26$~kJ: hematite plus liquid water is favored by only 0.13~kJ per mole of
FeOOH. This is well inside the $\pm1.4$~kJ\,mol$^{-1}$ uncertainty that \citet{majzlan2003b}
assign to the goethite--hematite free-energy difference; the two assemblages are essentially
balanced at ambient conditions. We carry this uncertainty ($\pm2.8$~kJ per reaction as written)
as bounding curves on the phase diagram.

\paragraph{Surface energy.} A crystal of finite size carries an excess free energy at its
surface in addition to the free energy of its bulk \citep{gibbs1878}. For a population of
particles of fixed size, this excess is $\gamma A$, where $\gamma$ is the surface energy
(J\,m$^{-2}$) and $A$ the total surface area. Both the bulk and surface terms scale with the
amount of material at fixed particle size, so the surface contribution can be written per mole,
$G = G_{\mathrm{bulk}} + \gamma A_m$, with $A_m$ the surface area per mole. $A_m$ grows as $1/d$,
so the term is negligible for coarse crystals and significant for nanoparticles
\citep{navrotsky2008}. For the reaction as written (one mole of hematite from two moles of
goethite),
\begin{equation}
  \Delta G_{\mathrm{surf}} = \gamma_{\mathrm{hem}} A_{m,\mathrm{hem}} - 2\,\gamma_{\mathrm{goe}} A_{m,\mathrm{goe}},
  \label{eq:gt_dGsurf}
\end{equation}
where $A_{m,\mathrm{hem}}$ is per mole of Fe$_2$O$_3$ and $A_{m,\mathrm{goe}}$ per mole of FeOOH.

A size must be chosen for each phase. We assume pseudomorphic dehydroxylation, in which hematite
inherits the external envelope of the parent goethite, so that both phases share the same
external surface area per mole of Fe. This is also the basis of the stability diagrams of
\citet{navrotsky2008}. For spherical goethite grains of diameter $d$ this shared area is
\begin{equation}
  A_{\mathrm{Fe}} = \frac{6\,V_{\mathrm{goe}}}{d}\quad\text{per mole of Fe}.
\end{equation}
Goethite has one Fe per formula unit, so $A_{m,\mathrm{goe}} = A_{\mathrm{Fe}}$. Hematite has two,
so $A_{m,\mathrm{hem}} = 2A_{\mathrm{Fe}}$. Substituting into Eq.~\ref{eq:gt_dGsurf},
\begin{equation}
  \Delta G_{\mathrm{surf}} = \gamma_{\mathrm{hem}}\,(2A_{\mathrm{Fe}}) - 2\,\gamma_{\mathrm{goe}}\,A_{\mathrm{Fe}}
  = 2A_{\mathrm{Fe}}\left(\gamma_{\mathrm{hem}} - \gamma_{\mathrm{goe}}\right).
  \label{eq:gt_equalarea}
\end{equation}
The factor of two counts the two moles of Fe that pass through the reaction. It multiplies both
terms because both phases carry the same area per Fe. For $d=10$~nm,
$A_{\mathrm{Fe}} = 6\times20.881\times10^{-6}/10^{-8} = 1.253\times10^{4}$~m$^2$\,mol$^{-1}$, and
with hydrated surface energies
$\Delta G_{\mathrm{surf}} = 2\times1.253\times10^{4}\times(0.75-0.60) = +3.76$~kJ.
The corresponding values are $+0.38$~kJ at 100~nm and $+0.04$~kJ at 1~$\mu$m. A positive
$\Delta G_{\mathrm{surf}}$ stabilizes goethite. The result depends on the choice of basis: at
equal diameter the same hydrated surface energies destabilize fine goethite
(Section~\ref{sec:gt_supp_basis}).

We use hydrated surface energies. They are referenced to liquid water, so they apply wherever the
surfaces are in equilibrium with water at or near saturation, which includes all of the
hydrothermal experiments. Anhydrous surface energies apply only to surfaces that have lost their
chemisorbed water, i.e.\ under vacuum.

\paragraph{Phase diagram.} The phase diagram (Fig.~\ref{fig:gt_PT}) is computed by solving
$\Delta G_r(T,P) = 0$ (Eqs.~\ref{eq:gt_dGr}--\ref{eq:gt_equalarea}) numerically for
$P = P_{\mathrm{H_2O}}$ at each temperature, using Brent's method on $\log_{10}P$ between
$10^{-12}$ and $10^{4}$~bar, with the solids at the fluid pressure. Where no root exists below the
melting curve of ice, the curve is truncated. Bulk goethite meets the liquid--vapor curve at
18~$^\circ$C, 100~nm goethite at 28~$^\circ$C and 10~nm goethite at 108~$^\circ$C. Above saturation,
water is liquid, the reaction volume
$\Delta V_r = V_{\mathrm{hem}} + V_{\mathrm{w}} - 2V_{\mathrm{goe}} \approx +6.6$~cm$^3$ is small, and the
boundaries become nearly vertical in $\log P$--$T$ space. At 25~$^\circ$C, bulk goethite requires
$P_{\mathrm{H_2O}}\approx10^{2.6}$~bar (compressed liquid). By contrast, 100~nm goethite is stable
above $10^{-1.52}$~bar (96\% relative humidity) and 10~nm goethite above $10^{-2.11}$~bar (24\%
relative humidity).

\begin{figure}
  \centering
  \includegraphics[width=0.7\linewidth]{goethite_hematite_Ph2o_T_hydrated}
  \caption{Goethite--hematite--water equilibrium in $\log P_{\mathrm{H_2O}}$--$T$ space for bulk
  goethite and for 100 and 10~nm grains with hydrated surfaces, compared at equal surface area
  per mole of Fe (Eq.~\ref{eq:gt_equalarea}). Goethite is stable above and to the left of each
  curve. Dash-dot curves: bulk boundary shifted by the $\pm1.4$~kJ\,mol$^{-1}$ FeOOH uncertainty of
  \citet{majzlan2003b}. Dotted: liquid--vapor curve of water. Thermodynamic data in
  Table~\ref{tab:gt_thermo_data}.}
  \label{fig:gt_PT}
\end{figure}

\paragraph{Driving force and the role of water.} The quantity that matters for kinetics is the
driving force for hematite formation, $\Delta\mu \equiv -\Delta G_r$ (per mole of Fe$_2$O$_3$).
Water controls $\Delta\mu$ through its own chemical potential, because it is a product of
reaction~\ref{eq:gt_rxn}. $G_{\mathrm{w}}$ in Eq.~\ref{eq:gt_dGr} is the chemical potential of
H$_2$O at the prevailing water fugacity. It depends on the partial pressure of water, not on the
total pressure. An inert gas, or any other mechanical load, acts only through the solid-volume
term $\Delta V_s(P-P_r)$, which contributes $-1.15$~J\,bar$^{-1}$ and slightly favors the denser
hematite. At 40~bar that is 46~J, three orders of magnitude smaller than the effect of 40~bar of
water. The water-buffered step-heating technique therefore stabilizes goethite primarily by
mass action on the bulk reaction. Surface hydration is a second-order correction that matters
only for nanometer grains (Eq.~\ref{eq:gt_equalarea}).

Table~\ref{tab:gt_dmu} gives the bulk driving forces at the experimental temperatures, evaluated
with IAPWS-95 water. At saturation, $\Delta\mu$ is 6--11~kJ, only 1.6--2.4 times $RT$. Under the
vacuum of our step-heating line ($P_{\mathrm{H_2O}}\approx10^{-8}$~atm) it is 76--106~kJ, about ten
times larger.

\begin{table}
  \centering
  \caption{Bulk driving force for hematite formation, $\Delta\mu = -\Delta G_r$ (kJ per mole Fe$_2$O$_3$), at water saturation and at $P_{\mathrm{H_2O}} = 10^{-8}$~atm.}
  \label{tab:gt_dmu}
  \begin{tabular}{cccccc}
    \hline
    $T$ ($^\circ$C) & $P_{\mathrm{sat}}$ (bar) & $\Delta\mu_{\mathrm{sat}}$ & $\Delta\mu_{\mathrm{sat}}/RT$ & $\Delta\mu_{\mathrm{vac}}$ & $(\Delta\mu_{\mathrm{vac}}/\Delta\mu_{\mathrm{sat}})^2$ \\
    \hline
    150 & 4.76  & 5.68  & 1.61 & 75.8  & 178 \\
    175 & 8.93  & 6.88  & 1.84 & 83.4  & 147 \\
    200 & 15.55 & 8.10  & 2.06 & 91.0  & 126 \\
    225 & 25.50 & 9.36  & 2.26 & 98.5  & 111 \\
    250 & 39.76 & 10.65 & 2.45 & 106.0 & 99  \\
    \hline
  \end{tabular}
\end{table}

% -----------------------------------------------------------------------------
\section{Vacuum step-heating is controlled by conversion to hematite}
\label{sec:gt_vacuum}
% -----------------------------------------------------------------------------

Two vacuum step-heating experiments were run on the same goethite. KAF used 1~h steps from
70 to 330~$^\circ$C, including isothermal repeats at 210, 230 and 240~$^\circ$C. Its extraction
line was open to a liquid-nitrogen-cooled charcoal trap, so evolved water did not accumulate,
and the water released in each step was measured manometrically. HM used a different schedule:
prograde to 238~$^\circ$C, retrograde to 180~$^\circ$C, then prograde again, with 0.5~h steps
above 238~$^\circ$C, to 498~$^\circ$C. Water could accumulate in the HM line. For each experiment, the cumulative
fraction of $^3$He released after step $i$ is
\begin{equation}
  F_i = \frac{\sum_{j\le i} {}^3\mathrm{He}_j}{\sum_{\mathrm{steps}} {}^3\mathrm{He} + {}^3\mathrm{He}_{\mathrm{fusion}}} .
\end{equation}

\paragraph{Kinetic model.} We describe conversion by a first-order rate law,
\begin{equation}
  \frac{d\alpha}{dt} = k(T)\,(1-\alpha),
  \qquad
  k(T) = A\exp\!\left(-\frac{E_a}{RT}\right),
  \label{eq:gt_F1}
\end{equation}
where $\alpha$ is the fraction of goethite converted. For a powder of many grains, first-order
kinetics is the limit in which each grain converts rapidly once a hematite nucleus forms, so
that the conversion rate is set by the nucleation rate per untransformed grain
\citep{mampel1940}. $k$ is then a per-grain nucleation rate. We assume $^3$He is released when
its host grain converts, except for a fraction $\phi$ that is inherited by the hematite
product:
\begin{equation}
  F = (1-\phi)\,\alpha .
  \label{eq:gt_Fphi}
\end{equation}
For a schedule of isothermal steps of duration $\Delta t_j$ at $T_j$, Eq.~\ref{eq:gt_F1}
integrates to
\begin{equation}
  \alpha_i = 1 - \exp\!\left[-A\sum_{j\le i} m_j\,\Delta t_j \exp\!\left(-\frac{E_a}{RT_j}\right)\right],
  \label{eq:gt_F1_cum}
\end{equation}
where $m_j = s$ for the two KAF steps run with the trap isolated (see below) and $m_j=1$
otherwise. Inverting Eq.~\ref{eq:gt_F1} step by step gives an apparent rate constant for every
step that is independent of any fitted parameter except $\phi$:
\begin{equation}
  k_i = -\frac{1}{\Delta t_i}\ln\frac{1-\alpha_i}{1-\alpha_{i-1}},
  \qquad \alpha_i = \frac{F_i}{1-\phi}.
  \label{eq:gt_stepk}
\end{equation}
For comparison, the diffusion parameter of each step is obtained with the
\citet{fechtig1966} inversion for a sphere,
\begin{equation}
  \frac{Dt}{a^2} =
  \begin{cases}
    \dfrac{1}{\pi^2}\left[2\pi - \dfrac{\pi^2}{3}F - 2\pi\sqrt{1-\dfrac{\pi F}{3}}\,\right], & F \le 0.85,\\[2ex]
    -\dfrac{1}{\pi^2}\ln\!\left[\dfrac{\pi^2}{6}(1-F)\right], & F > 0.85,
  \end{cases}
  \qquad
  \left(\frac{D}{a^2}\right)_i = \frac{(Dt/a^2)_i - (Dt/a^2)_{i-1}}{\Delta t_i}.
  \label{eq:gt_FK}
\end{equation}

\paragraph{Fitting window.} All steps are plotted, but fits use only KAF steps 5--23
(185--285~$^\circ$C) and HM steps 8--28 (199--338~$^\circ$C). KAF steps 1--4 and HM steps 2--7
(below $\sim$180~$^\circ$C) release a small ($<2\%$), separate low-temperature component
(Section~\ref{sec:gt_lowT}). Later steps sit on the post-conversion plateau, where
$1-\alpha\to0$ and Eq.~\ref{eq:gt_stepk} becomes controlled by the choice of $\phi$. The KAF
steps run with the trap isolated are excluded from the Arrhenius regression.

\paragraph{Results.} Fitting Eqs.~\ref{eq:gt_Fphi}--\ref{eq:gt_F1_cum} to the cumulative KAF
release by least squares gives $E_a = 122.9$~kJ\,mol$^{-1}$, $\ln A = 18.67$ ($A$ in s$^{-1}$),
$\phi = 0.112$ and $s = 0.46$, with an rms misfit in $F$ of 0.0085. The step rate constants of
both experiments (Eq.~\ref{eq:gt_stepk}) fall on a single Arrhenius line
(Fig.~\ref{fig:gt_arrhenius}a). This includes the HM retrograde steps, which a first-order
process requires: its rate depends only on $T$ and on the amount of goethite remaining. An
unweighted least-squares fit of $\ln k_i$ against $1/T_i$ for the 38 steps in the window gives
\begin{equation}
  E_a = 121.4 \pm 4.4~\mathrm{kJ\,mol^{-1}},
  \qquad
  \ln A = 18.08 \pm 1.05
  \label{eq:gt_line}
\end{equation}
(KAF alone $134.9\pm6.1$; HM alone $116.3\pm5.6$~kJ\,mol$^{-1}$), in agreement with the cumulative
fit. At 250~$^\circ$C this gives $k = 5.4\times10^{-5}$~s$^{-1}$, or 18\% conversion per hour. The
result does not depend on the fitting window: fitting all steps instead gives
$E_a = 122.7$~kJ\,mol$^{-1}$ and the same ranking of models (Section~\ref{sec:gt_supp_mech}).
Thermogravimetric studies of goethite dehydroxylation report somewhat higher apparent activation
energies, $154\pm15$ to $169\pm8$~kJ\,mol$^{-1}$, for a reaction controlled by a two-dimensional
phase boundary, with a different mechanism at early stages and low temperature \citep{goss1987}.
% TODO: compare conditions (heating in air/N2 vs vacuum; sample grain size) before finalizing.

In the diffusion representation (Fig.~\ref{fig:gt_arrhenius}b) the two experiments do not
collapse. HM's retrograde steps between 180 and 220~$^\circ$C plot about three natural-log units
above HM's own prograde steps at the same temperatures. Volume diffusion from grains that are
being progressively depleted cannot produce this, whether from a single domain or multiple
domains: retrograde steps must fall on or below the prograde trend. $D/a^2$ also turns over
above $\sim$275~$^\circ$C, where release stalls at the plateau.

\begin{figure}
  \centering
  \includegraphics[width=\linewidth]{arrhenius_reaction_vs_diffusion}
  \caption{Every step of the two vacuum experiments (KAF, HM) and the hydrothermal experiment
  in two representations. (a) First-order rate constants (Eq.~\ref{eq:gt_stepk},
  $\phi=0.112$); red line: fit to KAF and HM steps in the fitting window (filled symbols,
  Eq.~\ref{eq:gt_line}); open symbols: steps outside the window; purple circles: KAF steps with
  the LN$_2$ trap isolated. Steps with $F>1-\phi$ have no defined $k$. (b) \citet{fechtig1966}
  $D/a^2$ (Eq.~\ref{eq:gt_FK}). The vacuum experiments collapse in (a) but not in (b); the
  hydrothermal experiment shows the opposite behavior (Section~\ref{sec:gt_hydro}).}
  \label{fig:gt_arrhenius}
\end{figure}

\paragraph{Out-of-sample prediction.} The parameters fitted to KAF alone predict the independent
HM schedule with an rms misfit in $F$ of 0.017 over the fitting window
(Fig.~\ref{fig:gt_vac_fit}c,d). The prediction reproduces the main release and follows the
retrograde--prograde cycle within about a factor of two on individual steps. A single-domain
spherical diffusion model fits KAF worse (rms 0.046) and predicts HM worse (rms 0.053). It
also predicts $F\to1$ by $\sim$300~$^\circ$C, whereas both experiments stall at $F\approx0.89$.
Within a single schedule, fits alone cannot discriminate between kinetic mechanisms
(Section~\ref{sec:gt_supp_mech}). The second schedule is what gives the comparison its power.

\begin{figure}
  \centering
  \includegraphics[width=\linewidth]{vacuum_kinetics_KAF_to_HM}
  \caption{(a) Cumulative $^3$He release of KAF with fits of the first-order (F1),
  contracting-sphere (R3) and spherical-diffusion (D3) models; purple bands mark the steps with
  the LN$_2$ trap isolated; grey: temperature schedule. (b) $^3$He against water released during
  KAF, both normalized. (c) HM cumulative release and the predictions from the KAF-calibrated
  models, with no parameters refitted. (d) HM per-step release fractions.}
  \label{fig:gt_vac_fit}
\end{figure}

\paragraph{Whole-grain release.} % NOTE: remove this paragraph if 4He/3He is to be avoided entirely
During the main vacuum release (KAF, $0.04<F<0.88$), the step $^4$He/$^3$He ratio stays within
0.95--1.10 of the bulk ratio. Grains therefore release their integrated inventory as they
convert, rather than releasing a diffusive, rim-weighted component.

\paragraph{Retention through conversion.} In both experiments release stalls at
$F\approx0.89$. The fitted $\phi=0.112$ is the same for KAF and HM. In HM, the remaining
$\sim$11\% is released slowly between 340 and 500~$^\circ$C and on fusion. We interpret this
fraction as $^3$He inherited by the hematite product and subsequently released by the much
slower diffusion out of hematite. Single-domain goethite diffusion cannot produce a plateau.

\paragraph{Water release.} The water evolved in each KAF step tracks the $^3$He release closely
(Fig.~\ref{fig:gt_vac_fit}b). Water leads $^3$He slightly at 230~$^\circ$C and lags it by
$\sim$0.1 between 250 and 280~$^\circ$C. Fitting Eq.~\ref{eq:gt_F1_cum} to the water data alone
gives $E_a = 91.5$~kJ\,mol$^{-1}$ (rms 0.033). The larger misfit is consistent with
water adsorbed before dehydroxylation ($\sim$1\% of the total, released below 200~$^\circ$C). It
is also consistent with systematically low manometer readings on the first step of each
measurement day, and with normalizing the water to the 330~$^\circ$C total. % TODO: confirm the day-boundary artifact with the line operator
Water release thus provides an independent measure of conversion that is consistent with the
$^3$He record.

\paragraph{Trap-isolated steps.} In two KAF steps the trap was isolated and evolved water
accumulated in the line. At 230~$^\circ$C the step rate constant fell to $0.33$ times the mean of
its neighbors at the same temperature, and at 240~$^\circ$C to $0.22$ times. The rate recovered
fully when the trap was restored. Water and $^3$He were suppressed together. A diffusion
process would not respond to the partial pressure of water. We use these two steps in
Section~\ref{sec:gt_nucleation} to calibrate the nucleation model.

\paragraph{Low-temperature component.}
\label{sec:gt_lowT}
Below $\sim$180~$^\circ$C both experiments release $^3$He 0.5--3 orders of magnitude faster than
the extrapolated conversion rate (Fig.~\ref{fig:gt_delta}). This component amounts to less than
2\% of the $^3$He and is not conversion. We attribute it to diffusion from the least retentive
domains, and it is treated with the diffusion models of Section~\ref{sec:gt_hydro}.


% -----------------------------------------------------------------------------
\section{Hydrothermal release is diffusive}
\label{sec:gt_hydro}
% -----------------------------------------------------------------------------

The hydrothermal experiment heated an aliquot of the same goethite in distilled water, with
much of the sample submerged in the liquid. The schedule rose from 170 to 230~$^\circ$C, cycled
down to 175~$^\circ$C, rose to 250~$^\circ$C, cycled down to 225~$^\circ$C and back to 250~$^\circ$C,
where three 5~h steps preceded a 170.8~h step, and ended with a step at 240~$^\circ$C (442.8~h in
total). Only 14.8\% of the $^3$He was released. Under the vacuum kinetics of
Eq.~\ref{eq:gt_line}, the same schedule would have converted the sample completely. Two tests
show that the release follows diffusion rather than conversion.

\paragraph{Test 1: history dependence through temperature cycles.} In the first cycle,
$\ln(D/a^2)$ (Eq.~\ref{eq:gt_FK}) at 200 and 210~$^\circ$C is $-21.67$ and $-21.30$ on the first
prograde sweep. After the retrograde sweep to 175~$^\circ$C, the same temperatures give $-23.26$
and $-22.59$, lower by 1.3--1.6 natural-log units (Fig.~\ref{fig:gt_arrhenius}b). In the last
cycle, retrograde and prograde values at 230--245~$^\circ$C agree within 0.06--0.45 log units, and
the final prograde segment (225--250~$^\circ$C) is linear with $E_a = 157\pm3$~kJ\,mol$^{-1}$. This
is the signature of depleting fast, low-retentivity domains in the first cycle, with diffusion
dominated by more retentive domains thereafter. A first-order conversion rate depends only on
temperature and on the fraction converted, which is almost unchanged at $F\approx0.04$. It would
return retrograde steps to the prograde line, as the vacuum HM experiment does in the reaction
representation (Fig.~\ref{fig:gt_arrhenius}a).

\paragraph{Test 2: the 250~$^\circ$C holds.} The two processes project differently through a long
isothermal step. Diffusion slows as the diffusion front moves inward, whereas first-order
conversion proceeds at a nearly constant rate while $\alpha\ll1$. We calibrate both on the three
5~h steps at 250~$^\circ$C and project through the 170.8~h step. For first-order conversion,
$-\ln(1-\alpha)$ grows linearly in time at fixed $T$ (Eq.~\ref{eq:gt_F1}), so
\begin{equation}
  F(t) = (1-\phi)\left[1-(1-\alpha_0)\,e^{-kt}\right],
  \qquad \alpha_0 = \frac{F_0}{1-\phi},
\end{equation}
with $k$ fitted to the three calibration steps and $F_0 = 0.0793$ the release before the first of
them. For diffusion from a sphere, the Fourier number $Dt/a^2$ grows linearly in time at fixed
$T$:
\begin{equation}
  F(t) = 1 - \frac{6}{\pi^2}\sum_{n=1}^{\infty}\frac{1}{n^2}
  \exp\!\left[-n^2\pi^2\left(\left(\frac{Dt}{a^2}\right)_0 + \frac{D}{a^2}\,t\right)\right],
\end{equation}
where $(Dt/a^2)_0$ follows from $F_0$ by Eq.~\ref{eq:gt_FK} and $D/a^2$ is fitted to the
calibration steps. Both models reproduce the three calibration steps
(Fig.~\ref{fig:gt_hold}). Over the long step they diverge. Diffusion projects $F = 0.148$ against
an observed 0.146. First-order conversion projects 0.172, overpredicting the release in that step
by 44\%: $(0.172-0.087)/(0.146-0.087) = 1.44$. We also tested the contracting-geometry and Avrami
models of Section~\ref{sec:gt_supp_mech}, which fit the vacuum data less well. Calibrated the same
way, they project 0.174--0.209, and none changes the result.

\begin{figure}
  \centering
  \includegraphics[width=0.65\linewidth]{hydrothermal_diffusion_vs_conversion}
  \caption{The 250~$^\circ$C holds of the hydrothermal experiment. Black: observed cumulative
  $^3$He fraction. Curves: first-order conversion and spherical volume diffusion, each calibrated
  on the three 5~h steps (grey band) and projected through the 170.8~h step.}
  \label{fig:gt_hold}
\end{figure}

\paragraph{The reaction representation.} Plotted as first-order rate constants
(Fig.~\ref{fig:gt_arrhenius}a), the hydrothermal steps sit up to 2.8 orders of magnitude
(median 2.3) below the vacuum line. They show everything a first-order process forbids:
retrograde steps 2--4 natural-log units below prograde steps at the same temperature, and a
falling $k$ during the 250~$^\circ$C holds. In the diffusion representation
(Fig.~\ref{fig:gt_arrhenius}b) the same steps form a coherent, history-dependent array. The vacuum
experiments behave like conversion and the hydrothermal experiment behaves like diffusion.
Section~\ref{sec:gt_nucleation} shows that the conversion rate expected at water saturation is too
small, by tens of orders of magnitude, to contribute to the observed release.


% -----------------------------------------------------------------------------
\section{Why water suppresses conversion}
\label{sec:gt_nucleation}
% -----------------------------------------------------------------------------

\paragraph{Nucleation barrier.} In classical nucleation theory \citep{gibbs1878,volmer1926,becker1935},
reviewed by \citet{christian2002}, \citet{kashchiev2000} and, for mineral systems,
\citet{deyoreo2003} and \citet{lasaga1998}, a spherical hematite nucleus of radius $r$ has free
energy
\begin{equation}
  \Delta G(r) = -\frac{4}{3}\pi r^{3}\,\frac{\Delta\mu}{V_{\mathrm{hem}}} + 4\pi r^{2}\gamma,
\end{equation}
where $\Delta\mu$ is the driving force per mole of Fe$_2$O$_3$ (Table~\ref{tab:gt_dmu}),
$V_{\mathrm{hem}}$ the molar volume and $\gamma$ the energy of the interface surrounding the
nucleus. The first term is the free energy released by transforming the nucleus volume. The second
is the cost of creating its interface. Setting $d\Delta G/dr=0$ gives the barrier height
\begin{equation}
  \Delta G^{*} = \frac{16\pi\,\gamma^{3}\,V_{\mathrm{hem}}^{2}}{3\,\Delta\mu^{2}} .
  \label{eq:gt_barrier}
\end{equation}
Written per formula unit, $V_{\mathrm{hem}}\to V_{\mathrm{hem}}/N_A$ and $\Delta\mu\to\Delta\mu/N_A$, so
Avogadro's number cancels and Eq.~\ref{eq:gt_barrier} is directly the barrier per nucleus in
joules. In the first-order limit of Section~\ref{sec:gt_vacuum}, the rate constant is a nucleation
rate,
\begin{equation}
  k = A\exp\!\left(-\frac{E_a}{RT}\right)\exp\!\left(-\frac{\Delta G^{*}}{k_B T}\right),
  \qquad
  \frac{\Delta G^{*}}{k_B T} = \frac{C\gamma^{3}}{T\,\Delta\mu^{2}},
  \quad
  C \equiv \frac{16\pi V_{\mathrm{hem}}^{2}}{3k_B} = 1.112\times10^{15}~\mathrm{m^6\,K\,J^{-1}\,mol^{-2}} ,
  \label{eq:gt_knuc}
\end{equation}
where $A\exp(-E_a/RT)$ describes the atomic mobility and $\exp(-\Delta G^{*}/k_BT)$ the probability of
forming a critical nucleus. The barrier scales as $\gamma^3/\Delta\mu^2$. Between vacuum and water
saturation, $\Delta\mu$ falls tenfold and the barrier rises by the factor in the last column of
Table~\ref{tab:gt_dmu}: about 100 at 250~$^\circ$C and nearly 180 at 150~$^\circ$C.

\paragraph{An upper bound on the vacuum barrier.} The vacuum conversion rate is observed, so the
barrier that controls it cannot exceed what the fastest conceivable prefactor allows. In
transition-state theory \citep{eyring1935} no elementary step is faster than $k_BT/h$, which is
$1.09\times10^{13}$~s$^{-1}$ at 250~$^\circ$C. If every one of the $N$ formula units in a grain were a
nucleation site attempting at this frequency, the per-grain rate could not exceed
$N(k_BT/h)\exp(-\Delta G^{*}/k_BT)$. A 1~$\mu$m goethite sphere contains
$N = (\pi d^3/6)/(V_{\mathrm{goe}}/N_A) = 1.5\times10^{10}$ formula units. Requiring this ceiling to
reach the observed $k_{\mathrm{vac}} = 5.4\times10^{-5}$~s$^{-1}$ at 250~$^\circ$C (Eq.~\ref{eq:gt_line})
gives
\begin{equation}
  \frac{\Delta G^{*}_{\mathrm{vac}}}{k_BT} \le \ln\frac{N\,k_BT/h}{k_{\mathrm{vac}}}
  = \ln\frac{1.5\times10^{10}\times1.09\times10^{13}}{5.4\times10^{-5}} = 63 .
  \label{eq:gt_ceiling}
\end{equation}
A tenfold larger grain raises the bound by only $\ln 1000 = 7$. Real prefactors are smaller, so the
true barrier is lower still.

\paragraph{Nucleation must be internal.} The surface energies of \citet{navrotsky2008} fail this
bound by a wide margin. If the nucleus were bounded by a free, anhydrous hematite surface
($\gamma = 1.90$~J\,m$^{-2}$, $\Delta\mu = 106.0$~kJ at 250~$^\circ$C), Eq.~\ref{eq:gt_knuc} would give
$\Delta G^{*}/k_BT = 1.3\times10^{3}$. Nucleation on a goethite surface lowers this. For a nucleus
shaped as a spherical cap on a flat substrate, the barrier is reduced by the purely geometric
factor
$f(\theta) = (2+\cos\theta)(1-\cos\theta)^2/4 \le 1$, where the contact angle $\theta$ is set by the
balance of interfacial tensions at the triple junction (Young's equation),
$\cos\theta = (\gamma_{\mathrm{goe}} - \gamma_{\mathrm{hem|goe}})/\gamma_{\mathrm{hem}}$
\citep{volmer1926,turnbull1950}. Even with a hematite--goethite interfacial energy of zero,
$\cos\theta \le 0.91/1.90 = 0.48$, $f\ge0.17$ and $\Delta G^{*}/k_BT\ge2.2\times10^{2}$. No
non-negative interfacial energy brings nucleation on free surfaces within the bound of
Eq.~\ref{eq:gt_ceiling}. Hematite must therefore nucleate internally, where the only new
interface is between hematite and goethite.

This agrees with the known mechanism. Goethite dehydroxylates topotactically: the hematite
inherits the oxygen framework and orientation of the parent, and the pseudomorph develops pores
and twinned hematite domains \citep{limadefaria1963,watari1979,watari1983,cudennec2005}. Interfaces
of this kind are coherent or semicoherent. In the general classification of \citet{porter2009},
developed for metals, coherent interfaces have energies up to $\sim$0.2~J\,m$^{-2}$ and semicoherent
interfaces 0.2--0.5~J\,m$^{-2}$, compared with 0.5--1~J\,m$^{-2}$ for incoherent ones. For this system
specifically, density-functional calculations give 0.24~J\,m$^{-2}$ for the prismatic twin boundaries
between hematite domains formed by dehydrating goethite \citep{saito2016}. We are not aware of a
measured or calculated energy for the goethite--hematite interface itself. In what follows,
$\gamma$ is replaced by an effective interfacial energy $\gamma_{\mathrm{eff}}$, calibrated from the
data. It absorbs the geometry of internal nucleation and is a property of the sample.

\paragraph{Calibrating $\gamma_{\mathrm{eff}}$ from the trap-isolated steps.} Isolating the trap
raised $P_{\mathrm{H_2O}}$ in the line and lowered $\Delta\mu$, and the isothermal rate fell by a factor
$s$ (Section~\ref{sec:gt_vacuum}). With $A$ and $E_a$ unchanged, Eq.~\ref{eq:gt_knuc} gives
\begin{equation}
  \ln\frac{1}{s} = \frac{C\gamma_{\mathrm{eff}}^{3}}{T}\left(\frac{1}{\Delta\mu_{\mathrm{off}}^{2}} - \frac{1}{\Delta\mu_{\mathrm{on}}^{2}}\right)
  \quad\Longrightarrow\quad
  \gamma_{\mathrm{eff}} = \left[\frac{T\ln(1/s)}{C\left(\Delta\mu_{\mathrm{off}}^{-2} - \Delta\mu_{\mathrm{on}}^{-2}\right)}\right]^{1/3}.
  \label{eq:gt_gamma_cal}
\end{equation}
The driving forces are evaluated as follows.
\begin{enumerate}
  \item $\Delta\mu_{\mathrm{on}}$ and $\Delta\mu_{\mathrm{off}}$ are bulk driving forces,
  $-\Delta G_r$ from Eqs.~\ref{eq:gt_dGr}--\ref{eq:gt_Gw} with $\Delta G_{\mathrm{surf}} = 0$, at the
  step temperature and with the solids at vacuum.
  \item They differ only through the chemical potential of water: $\Delta\mu_{\mathrm{on}}$ is
  evaluated at $P_{\mathrm{H_2O}} = 10^{-8}$~atm and $\Delta\mu_{\mathrm{off}}$ at the water pressure
  reached with the trap isolated (below). At these pressures (at most 20~Pa), IAPWS-95 and the
  ideal gas agree, and
  $\Delta\mu_{\mathrm{on}} - \Delta\mu_{\mathrm{off}} = RT\ln(P_{\mathrm{off}}/P_{\mathrm{on}})$, which is
  32--41~kJ for the pressures in Table~\ref{tab:gt_gamma}.
  \item Surface terms are omitted in both. The barrier of an internal nucleus depends on the free
  energy per volume of material transformed, not on the external surface of the host grain. The
  grain or domain size of the natural sample is also not characterized. The external surface
  energy of the grain enters the overall thermodynamic stability (Section~\ref{sec:gt_thermo}),
  but we do not include it in the barrier. The sensitivity to this choice is small
  (Section~\ref{sec:gt_supp_barrier}).
\end{enumerate}
We bound the water pressure with the trap isolated by assuming that all evolved water was in the
vapor, distributed uniformly through an apparatus of volume $V_{\mathrm{app}}$ at 298~K:
\begin{equation}
  P_{\mathrm{H_2O}} = \frac{(m_{\mathrm{w}}/M_{\mathrm{w}})\,RT}{V_{\mathrm{app}}} .
\end{equation}
With $m_{\mathrm{w}} = 0.027$~mg at 230~$^\circ$C and 0.016~mg at 240~$^\circ$C ($1.50$ and
$0.89~\mu$mol), and $190\le V_{\mathrm{app}}\le600$~mL, Table~\ref{tab:gt_gamma} gives the pressures
and the resulting $\gamma_{\mathrm{eff}}$. The pressures are end-of-step values; we also evaluate half
of each as a time average for steady accumulation. Pressure enters only through $\ln P$ in
$\Delta\mu$ and then through a cube root, so a threefold uncertainty in volume changes
$\gamma_{\mathrm{eff}}$ by only about 10\%. The two steps agree within 15\%, and we adopt
\begin{equation}
  \gamma_{\mathrm{eff}} = 0.16\pm0.03~\mathrm{J\,m^{-2}} ,
\end{equation}
within the range expected for coherent and semicoherent interfaces. The corresponding barriers at
250~$^\circ$C are 0.5--1.3~$k_BT$ under vacuum and 52--129~$k_BT$ at saturation. Under vacuum the
barrier is negligible and the Arrhenius term of Eq.~\ref{eq:gt_line} sets the rate; at saturation
the barrier dominates.

\begin{table}
  \centering
  \caption{Water pressure with the trap isolated and the calibrated $\gamma_{\mathrm{eff}}$ (Eq.~\ref{eq:gt_gamma_cal}).}
  \label{tab:gt_gamma}
  \begin{tabular}{cccccc}
    \hline
    $T$ ($^\circ$C) & $s$ & $V_{\mathrm{app}}$ (mL) & $P_{\mathrm{H_2O}}$ (Pa) & $\gamma_{\mathrm{eff}}$, end-of-step $P$ & $\gamma_{\mathrm{eff}}$, $P/2$ \\
    \hline
    230 & 0.333 & 190 & 19.6 & 0.138 & 0.145 \\
    230 & 0.333 & 600 & 6.2  & 0.150 & 0.158 \\
    240 & 0.224 & 190 & 11.6 & 0.164 & 0.172 \\
    240 & 0.224 & 600 & 3.7  & 0.178 & 0.188 \\
    \hline
  \end{tabular}
\end{table}

\paragraph{Translation to water saturation.} For the same sample, with the same internal interface
and the same mobility term, only $\Delta\mu$ changes between vacuum and saturation:
\begin{equation}
  \frac{k_{\mathrm{sat}}}{k_{\mathrm{vac}}} = \exp\!\left[-\frac{C\gamma_{\mathrm{eff}}^{3}}{T}\left(\frac{1}{\Delta\mu_{\mathrm{sat}}^{2}} - \frac{1}{\Delta\mu_{\mathrm{vac}}^{2}}\right)\right],
  \label{eq:gt_translate}
\end{equation}
with both driving forces evaluated as in the calibration (bulk, IAPWS-95 water at $P_{\mathrm{sat}}$ and
at $10^{-8}$~atm). For $\gamma_{\mathrm{eff}} = 0.14$, 0.16 and 0.19~J\,m$^{-2}$, this ratio is
$7\times10^{-23}$, $9\times10^{-34}$ and $5\times10^{-56}$ at 250~$^\circ$C, and $5\times10^{-43}$,
$7\times10^{-64}$ and $2\times10^{-106}$ at 200~$^\circ$C.

How large a barrier would still allow detectable conversion follows from the same kinetics. Over
the 442.8~h ($1.59\times10^{6}$~s) hydrothermal experiment, converting even 1\% of the sample at
250~$^\circ$C requires $k_{\mathrm{sat}}t\gtrsim0.01$, i.e.
\begin{equation}
  \frac{\Delta G^{*}_{\mathrm{sat}}}{k_BT} \lesssim \ln\frac{k_{\mathrm{vac}}\,t}{0.01}
  = \ln\frac{5.4\times10^{-5}\times1.59\times10^{6}}{0.01} = 9 ,
\end{equation}
far below the calibrated 52--129~$k_BT$. Applying Eq.~\ref{eq:gt_translate} step by step to the
hydrothermal schedule, with $k_{\mathrm{vac}}$ from Eq.~\ref{eq:gt_line}, the predicted conversion is
\begin{equation}
  \alpha \approx \sum_i k_{\mathrm{vac}}(T_i)\,\frac{k_{\mathrm{sat}}}{k_{\mathrm{vac}}}(T_i)\,\Delta t_i
  = 10^{-20.6},\ 10^{-31.5},\ 10^{-53.7}
\end{equation}
for the three values of $\gamma_{\mathrm{eff}}$. Even the most permissive calibration,
$\gamma_{\mathrm{eff}} = 0.084$~J\,m$^{-2}$, gives $\alpha = 6\times10^{-4}$. That value would require
0.1~bar of water in the line, far above the volume constraint. If the observed release came from
conversion, it would imply $\alpha = F/(1-\phi) = 0.148/0.888 = 0.17$.

\paragraph{Residuals.} Figure~\ref{fig:gt_delta} shows the difference between each observed step
rate constant (Eq.~\ref{eq:gt_stepk}) and the rate predicted for its setting:
\begin{equation}
  \Delta\log_{10}k_i = \log_{10}k_i^{\mathrm{obs}} - \log_{10}k^{\mathrm{pred}}(T_i),
  \qquad
  k^{\mathrm{pred}} =
  \begin{cases}
    k_{\mathrm{vac}}(T) & \text{vacuum},\\
    k_{\mathrm{vac}}(T)\,\dfrac{k_{\mathrm{sat}}}{k_{\mathrm{vac}}}(T;\gamma_{\mathrm{eff}}) & \text{water-saturated}.
  \end{cases}
\end{equation}
The vacuum steps in the fitting window scatter about zero with an rms of 0.15 (KAF) and 0.21
(HM) orders of magnitude. The trap-isolated steps sit 0.30 and 0.55 orders below. The hydrothermal
steps sit 30--101 orders of magnitude above the saturated prediction for
$\gamma_{\mathrm{eff}}=0.16$~J\,m$^{-2}$, and 20--68 and 53--169 orders above for 0.14 and
0.19~J\,m$^{-2}$. Far more $^3$He is released in water than decomposition, as constrained by the vacuum
experiments, could supply. The hydrothermal release must be controlled by diffusion.

\begin{figure}
  \centering
  \includegraphics[width=\linewidth]{first_order_delta_lnk}
  \caption{First-order residuals, $\log_{10}k^{\mathrm{obs}} - \log_{10}k^{\mathrm{pred}}$, on a
  symmetric-log axis. Vacuum steps (KAF, HM) are referenced to the vacuum line
  (Eq.~\ref{eq:gt_line}); filled symbols are in the fitting window. Blue: hydrothermal steps
  referenced to the water-saturated prediction (Eq.~\ref{eq:gt_translate}) for
  $\gamma_{\mathrm{eff}} = 0.16$~J\,m$^{-2}$; bars span 0.14--0.19~J\,m$^{-2}$. Grey: the same hydrothermal
  steps referenced to the vacuum line without correction for water.}
  \label{fig:gt_delta}
\end{figure}

A rate law linear in the driving force cannot produce this contrast. Such a law has
$k\propto1-\exp(-\Delta\mu/RT)$, which at saturation is 0.80, 0.87 and 0.91 of the vacuum value at
150, 200 and 250~$^\circ$C. That predicts only a 10--20\% slowdown, far less than the trap-isolated
steps and hydrothermal data show. Suppression requires a barrier that grows as $\Delta\mu^{-2}$.

\paragraph{Assumptions.} This treatment is a consistent physical explanation of the data, not a
proof. It rests on four assumptions.
(i) The slowdown during the trap-isolated steps acts through $\Delta\mu$. If adsorbed water instead
slowed dehydroxylation by blocking reactive sites, these steps would not calibrate
$\gamma_{\mathrm{eff}}$.
(ii) The effective interface energy, $A$ and $E_a$ are the same in vacuum and in water. The
mechanism is internal nucleation in both settings, and only the chemical potential of water
changes.
(iii) The barrier depends on the bulk driving force. Including the surface terms of
Eq.~\ref{eq:gt_equalarea} for grain sizes down to 10~nm raises the predicted hydrothermal
conversion by about eight orders of magnitude, to $\alpha\approx10^{-23}$
(Section~\ref{sec:gt_supp_barrier}).
(iv) Only solid-state conversion is modeled. Conversion by dissolution of goethite and
precipitation of hematite from the fluid is not described by Eq.~\ref{eq:gt_translate}
(Section~\ref{sec:gt_pigment}).
The comparison is valid only between experimental conditions on the same sample. The calibrated
$\gamma_{\mathrm{eff}}$ absorbs the sample's microstructure and should not be transferred to other
goethites.


% -----------------------------------------------------------------------------
\section{Fluid-mediated conversion: synthetic pigment goethite}
\label{sec:gt_pigment}
% -----------------------------------------------------------------------------

Fluid-mediated conversion does occur in some materials. Acicular goethite blades
($\sim$1~$\mu$m long, aspect ratio $\sim$4), synthesized commercially as a yellow iron-oxide
pigment, were heated at 150--250~$^\circ$C under water-saturated conditions, both submerged in
liquid and in vapor. Some runs turned red. SEM showed that the blades had been replaced by
equant hematite plates $\sim$0.5~$\mu$m across, and both phases sometimes coexisted at
microscopic scales. Other runs at similar conditions did not react. The change of shape
rules out pseudomorphic dehydroxylation. It indicates dissolution of goethite and nucleation
and growth of hematite from the fluid, a pathway outside the solid-state model of
Section~\ref{sec:gt_nucleation}.

Grain size and bulk thermodynamics do not explain why these blades reacted when the natural
goethite did not. At 1~$\mu$m, the surface contribution to $\Delta G_r$ is only
$\sim$0.04~kJ, and the bulk driving force at saturation is the same 6--11~kJ for both materials.
The most likely difference is the fluid itself. Commercial iron-oxide pigments are commonly produced by
precipitation from iron-salt solutions. They can carry residual soluble salts, such as sulfate,
and may be surface-treated. Such contaminants would change the pH and ionic strength of the
water they are heated in, and with them the solubility of Fe(III). The dissolution--precipitation
pathway is limited by how much iron the fluid can carry, so a modest increase in Fe(III)
solubility could turn a negligible reaction into a visible one. Differences in residue loading
between vials would also explain why the outcome was not reproducible. The natural goethite, by
contrast, was heated in distilled water, where near-neutral Fe(III) solubility is very low.
% TODO: check the pigment product data sheet (grade, surface treatment, soluble-salt content, aqueous pH).

We conclude that water-saturated conditions suppress solid-state conversion of goethite, but
they do not by themselves exclude fluid-mediated conversion. Whether that pathway operates
depends on the solution chemistry. For the natural sample in distilled water,
Section~\ref{sec:gt_hydro} shows that the release is diffusive. For samples whose chemistry is
less well characterized, the absence of conversion should be verified directly, for example by
post-run X-ray diffraction or electron microscopy.


% =============================================================================
% SUPPLEMENT
% =============================================================================

% -----------------------------------------------------------------------------
\section{Kinetic mechanisms and the limits of a single heating schedule}
\label{sec:gt_supp_mech}
% -----------------------------------------------------------------------------

Solid-state rate laws take the form $d\alpha/dt = A\exp(-E_a/RT)\,h(\alpha)$. For any heating
schedule this integrates to
\begin{equation}
  g(\alpha) \equiv \int_0^{\alpha}\frac{d\alpha'}{h(\alpha')} = x,
  \qquad
  x = A\sum_j \Delta t_j\exp\!\left(-\frac{E_a}{RT_j}\right).
\end{equation}
We tested the mechanisms in Table~\ref{tab:gt_mech_forms}, including spherical diffusion on the
same footing. Under vacuum, the driving-force factor $1-\exp(-\Delta\mu/RT)$ is unity to better
than $10^{-9}$, so $x$ is simply an Arrhenius-weighted time. Diffusion and conversion then
differ only in the shape $\alpha(x)$ and in the fitted $E_a$.

\begin{table}
  \centering
  \caption{Kinetic models and their integrated forms.}
  \label{tab:gt_mech_forms}
  \begin{tabular}{llll}
    \hline
    Model & $h(\alpha)$ & $g(\alpha)$ & $\alpha(x)$ \\
    \hline
    R2 contracting cylinder & $2(1-\alpha)^{1/2}$ & $1-(1-\alpha)^{1/2}$ & $1-(1-x)^2$ \\
    R3 contracting sphere & $3(1-\alpha)^{2/3}$ & $1-(1-\alpha)^{1/3}$ & $1-(1-x)^3$ \\
    F1 first order & $1-\alpha$ & $-\ln(1-\alpha)$ & $1-e^{-x}$ \\
    A2, A3 Avrami ($n=2,3$) & $n(1-\alpha)[-\ln(1-\alpha)]^{1-1/n}$ & $[-\ln(1-\alpha)]^{1/n}$ & $1-e^{-x^n}$ \\
    D3 sphere diffusion & --- & --- & $1-\dfrac{6}{\pi^2}\displaystyle\sum_n \dfrac{e^{-n^2\pi^2 x}}{n^2}$ \\
    \hline
  \end{tabular}
\end{table}

\paragraph{Synthetic test.} We generated noiseless release curves from the R3 model
($E_a = 150$~kJ\,mol$^{-1}$) on two schedules: a monotonic series of 1~h steps from 150 to
390~$^\circ$C, and a cycled schedule with retrograde steps and long isothermal holds. We then
fitted every model with $E_a$ and $A$ free. Table~\ref{tab:gt_synth} shows that every model
except the Avrami forms fits within 1\% in $F$ by adjusting $E_a$. With realistic analytical
uncertainty, goodness of fit to a single schedule cannot identify the mechanism. The fitted
activation energy trades off against the shape of $\alpha(x)$. Discrimination requires a second,
independent schedule, an independent measure of $\alpha$, or both. The vacuum experiments supply
both (Section~\ref{sec:gt_vacuum}).

\begin{table}
  \centering
  \caption{Fits to synthetic R3 data ($E_a$ in kJ\,mol$^{-1}$; rms misfit in $F$).}
  \label{tab:gt_synth}
  \begin{tabular}{lcccc}
    \hline
    & \multicolumn{2}{c}{Monotonic schedule} & \multicolumn{2}{c}{Cycled schedule} \\
    Model & $E_a$ & rms & $E_a$ & rms \\
    \hline
    R2 & 140.0 & 0.0031 & 140.4 & 0.0028 \\
    R3 (true) & 150.0 & 0.0000 & 150.0 & 0.0000 \\
    F1 & 173.0 & 0.0074 & 177.4 & 0.0092 \\
    A2 & 81.7 & 0.0073 & 83.8 & 0.0077 \\
    A3 & 51.3 & 0.0073 & 53.8 & 0.0062 \\
    D3 & 296.2 & 0.0026 & 296.9 & 0.0025 \\
    \hline
  \end{tabular}
\end{table}

\paragraph{Fits to the vacuum data.} Table~\ref{tab:gt_mech_fits} gives the fit of each model to
the cumulative KAF $^3$He release within the fitting window. It also gives the misfit of the
resulting out-of-sample prediction of HM, and the fit to the KAF water data. F1 gives the best
fit and the best prediction. R3 and R2 are close behind, so we describe the mechanism as
nucleation- or interface-controlled rather than claim first order specifically. The Avrami
fits reach the lower bound of $E_a$ and fail on both data sets. Spherical diffusion fits KAF
five times worse than F1 and cannot produce the $F\approx0.89$ plateau.

The fitting window excludes the low-temperature steps, where release is dominated by diffusion from
goethite, and the plateau, where it is dominated by diffusion from hematite; the rms misfits are
computed over the same steps. Repeating the fits with every step included gives the same result:
F1 remains the best model ($E_a = 122.7$~kJ\,mol$^{-1}$, $\phi = 0.112$, rms 0.0073 for KAF and 0.026
for the HM prediction), followed by R3 (99.4~kJ\,mol$^{-1}$; 0.0097, 0.030) and R2 (87.9~kJ\,mol$^{-1}$;
0.0116, 0.033), with D3 at 175.9~kJ\,mol$^{-1}$ (0.063, 0.061).

\begin{table}
  \centering
  \caption{Model fits to the vacuum experiments (from \texttt{vacuum\_kinetics\_KAF\_to\_HM\_fits.csv}). $s$: rate factor for the trap-isolated steps.}
  \label{tab:gt_mech_fits}
  \begin{tabular}{lccccccc}
    \hline
    Model & $E_a$ & $\ln A$ & $\phi$ & $s$ & rms KAF & rms HM & $E_a$ (water) \\
    \hline
    R2 & 90.3 & 10.00 & 0.132 & 0.54 & 0.0133 & 0.0292 & 64.7 \\
    R3 & 101.7 & 12.38 & 0.127 & 0.50 & 0.0111 & 0.0234 & 73.4 \\
    F1 & 122.9 & 18.67 & 0.112 & 0.46 & 0.0085 & 0.0167 & 91.5 \\
    A2 & 25.6 & $-4.70$ & 0.000 & 0.20 & 0.0207 & 0.1814 & 20.2 \\
    A3 & 20.0$^{a}$ & $-5.89$ & 0.159 & 0.00 & 0.0446 & 0.2348 & 20.0$^{a}$ \\
    D3 & 183.7 & 29.50 & --- & 1.70 & 0.0455 & 0.0529 & 180.2 \\
    \hline
    \multicolumn{8}{l}{\footnotesize $E_a$ in kJ\,mol$^{-1}$, $A$ in s$^{-1}$. $^{a}$At the lower bound of the fit.}
  \end{tabular}
\end{table}


% -----------------------------------------------------------------------------
\section{Sensitivity of the size effect to the surface-area basis}
\label{sec:gt_supp_basis}
% -----------------------------------------------------------------------------

The sign of the surface stabilization depends on how the surface areas of the two phases are
related. At equal area per mole of Fe (Eq.~\ref{eq:gt_equalarea}), goethite is stabilized
whenever $\gamma_{\mathrm{hem}}>\gamma_{\mathrm{goe}}$. If both phases are instead spheres of the
same diameter $d$, then $A_i = 6V_i/d$ and
\begin{equation}
  \Delta G_{\mathrm{surf}} = \frac{6}{d}\left(\gamma_{\mathrm{hem}}V_{\mathrm{hem}} - 2\gamma_{\mathrm{goe}}V_{\mathrm{goe}}\right),
\end{equation}
which stabilizes goethite only if
\begin{equation}
  \frac{\gamma_{\mathrm{hem}}}{\gamma_{\mathrm{goe}}} > \frac{2V_{\mathrm{goe}}}{V_{\mathrm{hem}}} = \frac{2\times20.881}{30.274} = 1.379 .
\end{equation}
Hematite is denser and carries less surface per Fe at equal size. For hydrated surfaces,
$\gamma_{\mathrm{hem}}/\gamma_{\mathrm{goe}} = 0.75/0.60 = 1.25 < 1.379$, so fine goethite is
destabilized: at 10~nm,
$\Delta G_{\mathrm{surf}} = (6/10^{-8})(0.75\times30.274 - 2\times0.60\times20.881)\times10^{-6} = -1.41$~kJ.
For anhydrous surfaces, $1.90/0.91 = 2.09 > 1.379$ and $\Delta G_{\mathrm{surf}} = +11.7$~kJ at
10~nm. Table~\ref{tab:gt_basis} compares the two bases. The equal-area basis corresponds to
pseudomorphic conversion and underlies the main-text calculations. The equal-diameter basis is
more appropriate where hematite grows as separate crystals, as in fluid-mediated conversion
(Section~\ref{sec:gt_pigment}). There the surface term slightly favors hematite, but it is
negligible at micrometer sizes.

\begin{table}
  \centering
  \caption{$\Delta G_{\mathrm{surf}}$ (kJ per reaction as written) for the two surface-area bases.}
  \label{tab:gt_basis}
  \begin{tabular}{lcccc}
    \hline
    & \multicolumn{2}{c}{Hydrated} & \multicolumn{2}{c}{Anhydrous} \\
    $d$ & equal area & equal diameter & equal area & equal diameter \\
    \hline
    1~$\mu$m & $+0.04$ & $-0.01$ & $+0.25$ & $+0.12$ \\
    100~nm & $+0.38$ & $-0.14$ & $+2.48$ & $+1.17$ \\
    10~nm & $+3.76$ & $-1.41$ & $+24.8$ & $+11.7$ \\
    \hline
  \end{tabular}
\end{table}

\begin{figure}
  \centering
  \includegraphics[width=\linewidth]{size_basis_sensitivity}
  \caption{Goethite--hematite--water boundaries for hydrated surfaces at (a) equal surface area
  per mole of Fe and (b) equal particle diameter. The size effect changes sign between the two
  bases.}
  \label{fig:gt_basis}
\end{figure}


% -----------------------------------------------------------------------------
\section{Nucleation barrier ratio and the calibration of $\gamma_{\mathrm{eff}}$}
\label{sec:gt_supp_barrier}
% -----------------------------------------------------------------------------

From Eq.~\ref{eq:gt_barrier}, the ratio of barriers at water saturation and under vacuum at
the same temperature is
\begin{equation}
  \frac{\Delta G^{*}_{\mathrm{sat}}}{\Delta G^{*}_{\mathrm{vac}}}
  = \left(\frac{\Delta\mu_{\mathrm{vac}}}{\Delta\mu_{\mathrm{sat}}}\right)^{2}
    \left(\frac{\gamma_{\mathrm{sat}}}{\gamma_{\mathrm{vac}}}\right)^{3}.
\end{equation}
The prefactors, the temperature and the molar volume cancel. The first factor follows from the
thermodynamic data alone (Table~\ref{tab:gt_dmu}; 99 at 250~$^\circ$C for bulk goethite). The
second is 1 if nucleation occurs at the same internal interface in both settings. It is
$(0.75/1.90)^3 = 1/16$ in the extreme case that the relevant interface is a free hematite
surface, hydrated at saturation and anhydrous under vacuum. Figure~\ref{fig:gt_barrier_ratio}
shows the ratio between these bounds for bulk and 10~nm goethite. Even at the lower bound, the
saturated barrier is about six times the vacuum barrier at 250~$^\circ$C. Because the barrier
appears in an exponent, that is still a difference of many orders of magnitude in rate.

\begin{figure}
  \centering
  \includegraphics[width=0.7\linewidth]{nucleation_barrier_ratio_vs_T}
  \caption{Ratio of hematite nucleation barriers at water saturation and at
  $P_{\mathrm{H_2O}}=10^{-8}$~atm for bulk and 10~nm goethite. Solid: same interface in both
  settings; dashed: free hematite surfaces (hydrated vs anhydrous). Below the vertical dotted
  lines, goethite is stable at saturation and there is no driving force.}
  \label{fig:gt_barrier_ratio}
\end{figure}

\paragraph{Sensitivity to surface terms.} The calibration and translation use bulk driving
forces (Section~\ref{sec:gt_nucleation}). To test this choice, we repeat both with the surface
term of Eq.~\ref{eq:gt_equalarea} included in $\Delta\mu$. We use anhydrous surfaces for the vacuum
and trap-isolated steps and hydrated surfaces at saturation, the combination that most reduces the
vacuum driving force relative to the saturated one. Table~\ref{tab:gt_surf_sens} shows that the
conclusion is unchanged down to 10~nm grains.

\begin{table}
  \centering
  \caption{Effect of including surface terms in $\Delta\mu$ on the calibration and the predicted hydrothermal conversion. Driving forces in kJ per mole Fe$_2$O$_3$ at 250~$^\circ$C; $\gamma_{\mathrm{eff}}$ range over Table~\ref{tab:gt_gamma}.}
  \label{tab:gt_surf_sens}
  \begin{tabular}{lcccc}
    \hline
    Grain size & $\Delta\mu_{\mathrm{vac}}$ & $\Delta\mu_{\mathrm{sat}}$ & $\gamma_{\mathrm{eff}}$ (J\,m$^{-2}$) & $\log_{10}\alpha_{\mathrm{HT}}$ \\
    \hline
    Bulk & 106.0 & 10.65 & 0.138--0.188 & $-32.0$ \\
    100~nm & 103.5 & 10.27 & 0.133--0.182 & $-31.2$ \\
    10~nm & 81.2 & 6.89 & 0.090--0.131 & $-23.3$ \\
    \hline
    \multicolumn{5}{l}{\footnotesize $\alpha_{\mathrm{HT}}$: predicted conversion over the hydrothermal schedule, median $\gamma_{\mathrm{eff}}$.}
  \end{tabular}
\end{table}

Figure~\ref{fig:gt_translation}a shows $\gamma_{\mathrm{eff}}$ from Eq.~\ref{eq:gt_gamma_cal} as
a function of the assumed water pressure during the trap-isolated steps. It spans
0.08--0.20~J\,m$^{-2}$ over $10^{-5}$--$10^{-1}$~bar and 0.14--0.19~J\,m$^{-2}$ over the
volume-constrained range (Table~\ref{tab:gt_gamma}). Figure~\ref{fig:gt_translation}b shows the
predicted slowdown at saturation (Eq.~\ref{eq:gt_translate}) for the calibrated values and for
the most permissive value.

\begin{figure}
  \centering
  \includegraphics[width=\linewidth]{hydrothermal_translation}
  \caption{(a) $\gamma_{\mathrm{eff}}$ calibrated from the two trap-isolated KAF steps as a
  function of the assumed $P_{\mathrm{H_2O}}$; grey band: range allowed by the apparatus volume
  (190--600~mL) and the water released in each step. (b) Ratio of conversion rate constants at
  water saturation and under vacuum (Eq.~\ref{eq:gt_translate}); grey band: temperatures of the
  hydrothermal experiments.}
  \label{fig:gt_translation}
\end{figure}
